\section{引言}
\label{sec:intro}

再生数是判断传染病传播态势的基本指标，大于 1 意味着疫情扩张，小于 1 意味着收缩，常被用来评估干预效果和安排医疗资源。基于世代间隔与更新方程的估计方法已成为实时追踪的主要工具\citep{wallinga2004different, fraser2007estimating, cori2013new}，窗口选择、报告延迟与偏差修正等实际问题也有系统讨论\citep{thompson2019improved, gostic2020practical, nouvellet2018simple, steyn2025primer}。这些工作回答给定病例序列后怎样估计得更好；监测设计者面对的则是前一步问题：需要收集多少病例、覆盖多少机构才能达到所需精度。流感住院只占感染的 0.6\%--2\%\citep{zhou2012hospitalizations, tokars2018seasonal}，美国国家医疗安全网络（NHSN）各辖区每周住院数往往只有几十例，且填报制度在强制与自愿间切换\citep{cdc2026nhsn}。以这些数据为输入的集合预测已有成熟评价框架\citep{reich2019accuracy, cramer2022evaluation, mathis2024evaluation}，相对误差评分的性质也有专门研究\citep{bosse2023scoring}。设计者需要知道误差中有多大部分能靠扩大规模消除，提高报告率与扩大覆盖哪一条更划算，以及规模扩大到何处以后收益明显下降。

经典统计直觉给出了一个简单答案。在泊松更新模型下，再生数的 Fisher 信息随病例数线性增长，估计方差与病例数成反比\citep{parag2022quantifying}。超传播使继发感染数严重过离散，通常用负二项分布刻画\citep{lloydsmith2005superspreading, endo2020estimating, kucharski2020early}；它只改变方差的比例常数，不改变方差随病例数下降的速度。按照这一直觉，扩大规模总能换来同比例的精度提升，边际收益随规模平滑递减，不存在由传播过程本身决定的转折点。

真实监测中还有两类噪声使上述直觉需要修正。第一类是环境扰动。同一辖区、同一时段的感染者共享天气、接触模式与人群流动，学校开学、假期以及湿度与气温的变化使流感传播呈现季节性，各期再生数因而围绕均值波动\citep{dalziel2018urbanization, rasmussen2011inference}。生态学把这类波动称为环境随机性，以区别于个体层面的人口随机性。人口随机性的相对影响随种群规模增大而消失，环境随机性则不随规模减弱\citep{engen1998demographic, rohani2002determinism}。病例数增加能把瞬时再生数测得更准，却无法消除再生数本身的跨期波动。若监测的目的是把握一段时期的典型传播水平，或预测下一期传播，这部分波动就成为扩大规模无法消除的误差。它与 Eales 和 Riley（\citeyear{eales2024differences}）讨论的表观再生数不同：后者由一条病例序列与给定世代间隔反推，因卷积而滞后并平滑真实再生数，即使观测无误也存在；本文的两个估计目标分别是条件瞬时值与跨期均值，但在用滞后计数构造增长统计量时同样受到卷积影响。

第二类噪声来自报告过程。报告比例随就诊行为、检测能力与填报制度逐期变化，作用在整期报告数上，扩大覆盖人数同样无法消除。针对漏报与报告偏差，已有工作发展了修正方法与状态空间模型\citep{azmon2014estimating, stoner2019hierarchical, bracher2021marginal, breto2009time, bhatt2023semi, flaxman2020estimating}，EpiFilter 用相邻时期的信息平滑 $R_t$\citep{parag2021improved}。Parag 等\citep{parag2022quantifying, parag2022fundamental, parag2025asymmetric}量化了带噪发病曲线所含的信息，并给出识别疫情反弹所需时间的下限。这些工作关心给定数据能提供多少关于当期传播的信息；环境波动没有作为单独的方差成分出现，规模收益何时明显下降因而没有定量回答，报告波动与环境波动在计数中能否区分也没有现成答案。

本文以负二项传播与二项报告为基本模型，把监测目标分为当期瞬时再生数 $R_t$ 与跨期平均再生数 $\Rbar$。对估计目标1，$\widehat R_t=C_t/m$ 达到 Cramér--Rao 下界，方差按 $1/m$ 趋于零，离散参数趋于无穷时退化为泊松结果。对估计目标2，条件无偏估计的方差为环境方差 $v_R$ 与规模项 $A/m$ 之和；报告率逐期波动时再加入报告异质性方差 $v_H$。规模收益开始明显减小的位置由方差平衡规模 $\mx=A/v_R$ 给出，监测设计由此归结为规模、期数与报告口径三个杠杆。这些结论来自 Cramér--Rao 不等式、全方差公式、Fisher 信息的数据处理不等式与单调似然比检验\citep{casella2002statistical, lehmann2005testing, zamir1998proof}，预算分配与两阶段抽样的费用分配同构\citep{cochran1977sampling}。报告率波动下，提高报告率与扩大覆盖的优劣存在反转阈值，而两条途径都不改变 $v_H$；仅凭计数的前两阶矩无法区分 $v_R$ 与 $v_H$，但报告噪声使预测误差的截距按 $F(w,\psi)$ 随滞后权重变化，这一签名可以用不同分母检验。规模足够大以后，$v_R$ 本身的估计精度只由有效期数 $T_{\mathrm{eff}}\approx T(1-\phi^2)/(1+\phi^2)$ 决定，季层面方差的相对标准误不低于 $\sqrt{2/S}$。

应用部分使用 NHSN 2020--2026 年的 16,218 个州周流感住院周报，其中 10,303 个州周满足规模条件，流行初期样本为 1,267 个州周。单步预测误差随规模按 $a+b/D_t$ 下降，截距为 0.0991（95\% 区间 0.0837--0.1165），斜率为 2.243（1.494--3.007），经验比值 $b/a=22.6$（13.5--34.9），与情景参数给出的平衡规模处于同一量级；不含底板的均值形状明显更差。截距高于环境方差情景。用各州每周填报医院比例调整计数后，截距变化为 $-0.0004$（90\% BCa 区间 $-0.0012$ 至 $0.0005$），可由填报比例观测到的报告成分不足截距超出量的 1\%，填报比例以外的报告离散不在这一上界之内。截距比随滞后权重变化的方向在流感、COVID-19 与 RSV 的 7 个病原与窗口组合中一致，但入院延迟在环境各期独立时也会产生这一方向；按联合判据，4 个组合部分复现，3 个只呈方向一致，没有组合复现。季间截距差异显著，季层面方差只能靠积累流行季确定。疫情预警与趋势估计对监测规模的需求不同，宜分开设计。
